\section{Exponential families}\label{sec:expFamily}

It is well known that exponential families are specific maximum entropy distributions.
However, not any maximum entropy distribution is an exponential family member.
To prepare for the characterization of generic maximum entropy distributions, we will in this section represent exponential family members as \ComputationActivationNetwork{}s.
We further show that a maximum entropy distribution is a exponential family member, if and only if its mean parameter is in the relative interior of the mean polytope.

\subsection{Representation as elementary \ComputationActivationNetworks{}}

Let us first define exponential families as classes of probability distributions and then derive their representation by elementary \ComputationActivationNetworks{}.

\begin{definition}[Exponential Family]
    \label{def:expFamily}
    Given a statistic function
    \begin{align*}
        \sstat \defcols \facstates \rightarrow \parspace
    \end{align*}
    and a base measure $\basemeasurewith$ with $\contraction{\basemeasure}\neq0$, the set $\expfamily=\{\expdist\wcols\canparamwithin\}\subset\distmeasuredby{\basemeasure}$ of probability distributions
    \begin{align*}
        \expdistat{\shortcatvariables}
        = \frac{
            \expof{\contractionof{\sencsstatat{\shortcatvariables,\selvariable},\canparamwith}{\shortcatvariables}}
        }{
            \contraction{\expof{\contractionof{\sencsstatat{\shortcatvariables,\selvariable},\canparamwith}{\shortcatvariables}},\basemeasurewith}
        }
    \end{align*}
    is called the exponential family to $\sstat$.
\end{definition}

% Relation to standard notation and information geometry
To connect with standard notation of exponential families we notice that the coordinate of any distribution $\expdistat{\shortcatvariables}$ at some state $\shortcatindicesin$ is
\begin{align*}
    \expdistat{\indexedshortcatvariables}
    = \expof{\sum_{\selindexin}\weightparamat{\indexedselvariable}\cdot\sstatcoordinateofat{\selindex}{\shortcatindices} - \psi_{\sstat,\basemeasure}(\weightparam)}
\end{align*}
where
\begin{align*}
    \psi_{\sstat,\basemeasure}(\weightparam)
    = \lnof{\contraction{\expof{\contractionof{\sencsstatat{\shortcatvariables,\selvariable},\canparamwith}{\shortcatvariables}},\basemeasurewith}}
\end{align*}
is the cumulant generating function.
The function $\psi_{\sstat,\basemeasure}(\weightparam)$ and its Legendre transform $\psi^*_{\sstat,\basemeasure}(\meanparam)$ define the dually flat Riemannian structure on the exponential family manifold \cite{amari_information_2016}.
%
We now continue with the notation of \defref{def:expFamily} to find representations of exponential families by elementary \ComputationActivationNetworks{}.


\begin{theorem}[Exponential Families are elementary \ComputationActivationNetworks{}]
    \label{the:expFamilyTensorRep}
    Given any base measure $\basemeasure$ and a sufficient statistic $\sstat$ we enumerate for each coordinate $\selindexin$ the image $\imageof{\sstatcoordinateof{\selindex}}$ by a variable $\headvariableof{\selindex}$ taking values in $[\cardof{\imageof{\sstatcoordinateof{\selindex}}}]$, given an interpretation map
    \begin{align*}
        \indexinterpretationof{\selindex} \defcols
        [\cardof{\imageof{\sstatcoordinateof{\selindex}}}] \rightarrow \imageof{\sstatcoordinateof{\selindex}} \, .
    \end{align*}
    For any canonical parameter vector $\canparamwithin$ we build the activation cores $\softactlegwith$ for each coordinate $\headindexof{\selindex}\in[\cardof{\imageof{\sstatcoordinateof{\selindex}}}]$ by
    \begin{align*}
        \softactleg\left[\indexedheadvariableof{\selindex}\right]
        = \expof{\canparamat{\indexedselvariable} \cdot \indexinterpretationofat{\selindex}{\headindexof{\selindex}} } \,
    \end{align*}
    and have (see \figref{fig:expdistUnaryRealizable})
    \begin{align*}
        \expdistat{\shortcatvariables}
        = \frac{
            \contractionof{\{\bencsstatwith\}\cup\{\softactlegwith \wcols \selindexin\}}{\shortcatvariables}
        }{
            \contraction{\{\bencsstatwith\}\cup\{\softactlegwith \wcols \selindexin\}\cup\{\basemeasurewith\}}
        } \, .
    \end{align*}
\end{theorem}
\begin{proof} %% SHORTENED THE PROOF IN THE REPORT
    For each $\shortcatindicesin$ we have
    \begin{align*}
        &\contractionof{\{\bencsstatwith\}\cup\{\softactlegwith \wcols \selindexin\}}{\indexedshortcatvariables} \\
        &\quad = \prod_{\selindexin} \contraction{\bencodingofat{\sstatcoordinateof{\selindex}}{\headvariableof{\selindex},\indexedshortcatvariables},\softactlegwith} \\
        &\quad = \prod_{\selindexin} \expof{\canparamat{\indexedselvariable}\cdot\indexinterpretationofat{\selindex}{\sstatcoordinateofat{\selindex}{\shortcatindices}}} \\
        &\quad = \expof{\sum_{\selindexin}\canparamat{\indexedselvariable}\cdot\indexinterpretationofat{\selindex}{\sstatcoordinateofat{\selindex}{\shortcatindices}}} \\
        &\quad = \expof{\contractionof{\sencsstatat{\shortcatvariables,\selvariable},\canparamwith}{\indexedshortcatvariables}} \, .
    \end{align*}
    Therefore we have
    \begin{align*}
        \contractionof{\{\bencsstatwith\}\cup\{\softactlegwith \wcols \selindexin\}}{\shortcatvariables}
        = \expof{\contractionof{\sencsstatat{\shortcatvariables,\selvariable},\canparamwith}{\shortcatvariables}} \, .
    \end{align*}
    The claim follows, since this implies that also the contraction of both sides with $\basemeasurewith$ is equivalent.
\end{proof}

\begin{figure}[t]
    \begin{center}
        \input{tikz/exp-family/expdist_unary_realizable}
    \end{center}
    \caption{Representation of a member in an exponential family by a \ComputationActivationNetwork{} with elementary activation.
    Here $\partitionfunction\in\rr$ is a normalization constant.}\label{fig:expdistUnaryRealizable}
\end{figure}

We will use the following well known property (see e.g. \cite{brown_fundamentals_1987}), that there is a one-to-one map between the canonical parameters and the mean parameters.

\begin{lemma}
    \label{lem:interiorRepExpFamily}
    The set of mean parameters of the members of an exponential family is the relative interior of the mean polytope.
    %% NEEDED?
%    For $\canparam,\seccanparam\in\parspace$ with
%    \begin{align*}
%        %\uniquantwrtof{\selindexin}{
%            \contractionof{\expdistofat{\sstat,\canparam,\basemeasure}{\shortcatvariables},\sencsstatwith,\basemeasurewith}{\selvariable}
%            = \contractionof{\expdistofat{\sstat,\seccanparam,\basemeasure}{\shortcatvariables},\sencsstatwith,\basemeasurewith}{\selvariable}
%        %}
%    \end{align*}
%    we furthermore have $\expdistofat{\sstat,\canparam,\basemeasure}{\shortcatvariables}=\expdistofat{\sstat,\seccanparam,\basemeasure}{\shortcatvariables}$.
\end{lemma}
\begin{proof}
    See The~3.3 in \cite{wainwright_graphical_2008}.
\end{proof}

Based on this property we define the forward and backward mappings to an exponential family.

\begin{definition}\label{def:forwardBackwardMapping}
    The forward map of an exponential family is the map
    \begin{align*}
        \forwardmap  \defcols \parspace \rightarrow \sbinteriorof{\genmeanset}
    \end{align*}
    defined as
    \begin{align*}
        \forwardmapof{\canparam}
        = \contractionof{\expdistofat{\sstat,\canparam,\basemeasure}{\shortcatvariables},\sencsstatwith,\basemeasurewith}{\selvariable}
    \end{align*}
    Any map $\backwardmap: \sbinteriorof{\genmeanset}\rightarrow \parspace$ with $\forwardmapof{\backwardmapof{\meanparam}}=\meanparam$ for all $\meanparam\in\sbinteriorof{\genmeanset}$ is called a backward map.
\end{definition}

\subsection{Maximum entropy distributions}

One universal property of exponential families amongst others (see \cite{murphy_probabilistic_2023}) is that they maximize the entropy among those with same mean vector.

\begin{lemma}\label{lem:expFamilyMaxEnt}
    Any member of an exponential family $\expfamilyof{\sstat,\basemeasure}$ is a maximum entropy distribution with respect to the moment constraint on $\sstat$ and the base measure $\basemeasure$.
\end{lemma}
\begin{proof}
    Let $\expdistof{\sstat,\canparam,\basemeasure}$ be a exponential family member with canonical parameter $\canparam$ and let $\meanparamwith$ be the corresponding mean parameter.
    Let $\secprobat{\shortcatvariables}$ be any feasible distribution for the maximum entropy problem with respect to the moment constraint on $\sstat$ by $\meanparamwith$ and the base measure $\basemeasure$.
    We also have $\contractionof{\secprobat{\shortcatvariables},\sencsstatwith,\basemeasurewith}{\selvariable}=\meanparamat{\selvariable}$ and thus
    \begin{align*}
        \centropyofwrt{\secprobtensor}{\expdistof{\sstat,\canparam,\basemeasure}}{\basemeasure}
        &= -\contraction{\secprobat{\shortcatvariables},\lnof{\expdistofat{\sstat,\canparam,\basemeasure}{\shortcatvariables}},\basemeasurewith} \\
        &= -\contraction{\secprobat{\shortcatvariables},\sencsstatwith,\canparamwith,\basemeasurewith} + \cumfunctionof{\canparam} \\
        &= - \contraction{\canparamwith,\meanparamwith} + \cumfunctionof{\canparam} \\
        &= \sentropyofwrt{\expdistof{\sstat,\canparam,\basemeasure}}{\basemeasure} \, .
    \end{align*}
    With the Gibbs inequality we have if $\secprobtensor\neq\expdistof{\sstat,\canparam,\basemeasure}$
    \begin{align*}
        \sentropyofwrt{\expdistof{\sstat,\canparam,\basemeasure}}{\basemeasure}  - \sentropyofwrt{\secprobtensor}{\basemeasure}
        = \centropyofwrt{\secprobtensor}{\expdistof{\sstat,\canparam,\basemeasure}}{\basemeasure}  - \sentropyofwrt{\secprobtensor}{\basemeasure}  > 0 \,
    \end{align*}
    and thus $\sentropyofwrt{\secprobtensor}{\basemeasure}<\sentropyofwrt{\expdistof{\sstat,\canparam,\basemeasure}}{\basemeasure}$.
    Therefore, any feasible $\secprobtensor$ different from $\expdistof{\sstat,\canparam,\basemeasure}$ have less entropy than $\expdistof{\sstat,\canparam,\basemeasure}$ and $\expdistof{\sstat,\canparam,\basemeasure}$ is a maximum entropy distribution.
\end{proof}

Together with \lemref{lem:interiorRepExpFamily} stating the existence of an exponential family member for any mean parameter in the relative interior of the mean polytope we arrive at the following theorem.

\begin{theorem}
    \label{the:maxEntropyInterior}
    If and only if $\meanparamwith$ is in the relative interior of $\genmeanset$, then the unique solution of the maximum entropy problem is the distribution
    \begin{align*}
        \expdistofat{\sstat,\weightparam,\basemeasure}{\shortcatvariables}\in\expfamilyof{\sstat,\basemeasure}
    \end{align*}
    where for some backward map $\backwardmapwrt{\sstat,\basemeasure}$
    \begin{align*}
        \weightparam
        = \backwardmapwrtof{\sstat,\basemeasure}{\meanparam} \, .
    \end{align*}
\end{theorem}
\begin{proof}
    By \lemref{lem:interiorRepExpFamily} if
    \begin{align*}
        \meanparamwith \in \sbinteriorof{\genmeanset}  \, ,
    \end{align*}
    there is at least one weight parameter $\weightparamwith\in\parspace$ with
    \begin{align*}
        \contractionof{\expdistofat{\sstat,\weightparam,\basemeasure}{\shortcatvariables},\sencsstatwith,\basemeasurewith}{\selvariable}
        =\meanparamat{\selvariable}
    \end{align*}
    and $\backwardmapwrtof{\sstat,\basemeasure}{\meanparam}$ with respect to any backward map $\backwardmapwrt{\sstat,\basemeasure}$ is such one.
    With \lemref{lem:expFamilyMaxEnt} we conclude that $\expdistof{(\sstat,\backwardmapwrtof{\sstat,\basemeasure}{\meanparam},\basemeasure)}$ is the maximum entropy distribution with respect to the moment constraint on $\sstat$ by $\meanparamwith$ and the base measure $\basemeasure$.
\end{proof}

% Minimality comments
While the maximum entropy distribution is unique, the weight vector $\weightparam$ representing the maximum entropy distribution is not always unique.
This is reflected by multiple possible choices of backward mappings.
Only if the exponential family is minimal (see \cite{wainwright_graphical_2008}) then the weight vector is unique.
However, we did not restrict to such situations, since exponential families on refined base measures are not minimal.